\slajdid{09_1-s012}
\begin{frame}[label=09_1-s012]
\gusttekst
\setlength{\abovedisplayskip}{3pt}
\setlength{\belowdisplayskip}{3pt}
\setlength{\abovedisplayshortskip}{2pt}
\setlength{\belowdisplayshortskip}{2pt}
Analiza invertora. Od njega lako napraviti složena a da ostanu „iste“ karakteristike.\\
Prvi - \alert{nMOS invertor sa pasivnim opterećenjem}

\begin{columns}[c,onlytextwidth]
\begin{column}{.46\textwidth}\centering
\mbox{\begin{circuitikz}[DEoznaka,baseline=(current bounding box.center),circuitikz/bipoles/length=8mm]
\node[nigfete,scale=.7] (t) at (0,0) {};
\draw[DEvod] (t.D) -- (0,.6) to[R,l=\(R_L\)] (0,1.4) -- (0,1.55);
\draw[DEvod] (-.2,1.55)--(.2,1.55) node[above left] {\(V_{DD}\)};
\draw[DEvod] (t.S)--++(0,-.25) node[ground] {};
\draw[DEvod] (t.G)--++(-.3,0) node[ocirc] {} node[left] {\(V_I\)};
\draw[DEvod] (0,.6)--(.3,.6) node[ocirc] {} node[right] {\(V_O\)};
\node[right] at (.15,0) {\(T_D\)};
\end{circuitikz}
\hspace{1mm}%
% Kvalitativna kriva iz originala: nema uvedene numeričke skale.
\begin{tikzpicture}[DEoznaka,baseline=(current bounding box.center)]
\begin{axis}[width=25mm,height=20mm,scale only axis,
 axis lines=middle,xmin=-.08,xmax=1.22,ymin=-.06,ymax=1.2,
 ticks=none,clip=false,xlabel=\(V_I\),ylabel=\(V_O\),
 every axis x label/.style={at={(axis description cs:1,0)},anchor=north},
 every axis y label/.style={at={(axis description cs:0,1)},anchor=south},
 axis line style={DEvod,-{Latex[length=1.5mm]}}]
\addplot[DEvod,smooth] coordinates {(0,1)(.13,1)(.24,.98)(.35,.85)(.43,.55)(.52,.22)(.60,.12)(.85,.085)(1,.06)};
\draw[dashed] (axis cs:.13,0)--(axis cs:.13,1.1);
\draw[dashed] (axis cs:.59,0)--(axis cs:.59,1.1);
\draw[dashed] (axis cs:1,0)--(axis cs:1,1.1);
\node[above] at (axis cs:.075,1.02) {1.};
\node[above] at (axis cs:.34,1.02) {2.};
\node[above] at (axis cs:.80,1.02) {3.};
\node[left] at (axis cs:0,1) {\(V_{DD}\)};
\node[below] at (axis cs:1,0) {\(V_{DD}\)};
\end{axis}
\end{tikzpicture}

}
\end{column}
\begin{column}{.53\textwidth}\gusttekst\centering
U prvoj oblasti
\[
V_I=V_{GS,D}<V_{Tn}
\]
i tranzistor ne \(T_D\) (D – drive) vodi. Napon na izlazu
\[
\begin{aligned}
V_O&=V_{DD}-R_LI_{RL}\\
&=V_{DD}-R_LI_{Dn,D}=V_{DD}
\end{aligned}
\]
\end{column}
\end{columns}

Kada ulazni napon postane jednak \(V_{Tn}\) tranzistor počinje da vodi ali sa malom strujom. Zbog toga ne dolazi do velikog pada napona na otporniku \(R_L\) (L – load) odnosno napon između drejna i sorsa tranzistora će i dalje biti visok.
\[
V_{DSn}>V_{DSnsat}
\]
\[
V_{DSnsat}=
\frac{(\alert{V_{GSn}}-V_{Tn})L_nE_{Cn}}{L_nE_{Cn}+(\alert{V_{GSn}}-V_{Tn})}
=\frac{(\alert{V_{Tn}+\varepsilon}-V_{Tn})L_nE_{Cn}}{L_nE_{Cn}+(\alert{V_{Tn}+\varepsilon}-V_{Tn})}
\approx\alert{\varepsilon}
\]
\[
V_O=V_{DD}-R_LI_{RL}=V_{DD}-R_LI_{Dn,D}<V_{DD}
\]
\end{frame}
